\section{Data and Preprocessing}

The data were obtained from the Alzheimer's Disease Neuroimaging
Initiative (ADNI)~\cite{ADNI}. Two datasets were used: a primary dataset
(\textbf{ADNI-B}) and a longitudinal dataset (\textbf{ADNI-Long}).
Both contain resting-state functional magnetic resonance imaging
(rs-fMRI), regional amyloid-$\beta$ (A$\beta$) and tau measurements,
and a clinical diagnosis for each session. Diagnoses were grouped as
healthy controls (HC), mild cognitive impairment (MCI), and Alzheimer's
disease (AD).

ADNI-B contains 152 ADNI3 participants, with one session per
participant. It was used for model development and evaluation.
ADNI-Long contains data from ADNI1, ADNIGO, ADNI2, and ADNI3. Before
longitudinal selection, it included 141 participants and 281 sessions.
Its repeated sessions allow changes in clinical stage to be studied
across visits.

Both datasets were provided in preprocessed and parcellated form.
Neuroimaging preprocessing, denoising, spatial normalisation, temporal
filtering, and brain parcellation were therefore not performed in this
work. The datasets were processed independently using the same
procedure, so matching sessions are not necessarily numerically
identical. Sections~\ref{sec:neuroimaging_preprocessing}
and~\ref{sec:dataset_consistency} describe the upstream procedure and
the comparison between datasets, respectively.


\subsection{Sample Selection and Cohort Characteristics}
\label{sec:sample_selection}

ADNI-B was the cross-sectional cohort for training, validation, and
testing. Its 152 samples represent different participants: 93 HC, 33
MCI, and 26 AD.

ADNI-Long was used to evaluate changes associated with clinical
progression. Participants with unchanged diagnoses across all available
visits were excluded because the analysis focuses on changes between
clinical stages. Transitions to an earlier stage were also excluded.
The selected cohort therefore represents forward progression between
HC, MCI, and AD rather than all trajectories in ADNI-Long.

Participants with three sessions and different diagnoses were split
into pairs representing clinical progression. The construction of these
pairs is described in Section~\ref{sec:longitudinal_construction}.

Participants shared by the selected ADNI-Long cohort and ADNI-B were
identified. To avoid subject-level information leakage, their ADNI-B
samples were assigned only to the test set. They were therefore not
included in the training or validation data.

Table~\ref{tab:dataset_summary} summarises ADNI-B, the available
ADNI-Long dataset, and the selected longitudinal cohort.

\begin{table}[ht]
    \centering
    \caption{Summary of the available and selected datasets. Diagnostic
    counts for ADNI-Long are reported at the session level.}
    \label{tab:dataset_summary}
    \begin{tabular}{lccc}
        \hline
        & ADNI-B & ADNI-Long Available & ADNI-Long Selected \\
        \hline
        ADNI phases & ADNI3 & ADNI1/GO/2/3 & ADNI1/2/3 \\
        Unique subjects & 152 & 141 & 22 \\
        Sessions & 152 & 281 & 48 \\
        HC sessions & 93 & 175 & 15 \\
        MCI sessions & 33 & 76 & 23 \\
        AD sessions & 26 & 30 & 10 \\
        \hline
    \end{tabular}
\end{table}

Table~\ref{tab:demographics} shows the demographic characteristics of
the final cohorts. Mean age was similar across groups in ADNI-B
(73.26--75.88 years) and ADNI-Long (73.87--74.67 years). Sex was also
relatively balanced, although the AD group in ADNI-B contained more
males (61.5\%).

ADNI-Long characteristics are reported at the session level. A
participant may therefore contribute to more than one clinical group as
their diagnosis changes across visits.

\begin{table}[ht]
    \centering
    \caption{Demographic characteristics of the final ADNI-B and selected
    ADNI-Long cohorts. Age is reported as mean $\pm$ standard deviation.
    ADNI-Long characteristics are reported at the session level.}
    \label{tab:demographics}
    \begin{tabular}{llcccc}
        \hline
        Dataset & Group & N & Age (years) & Female, n (\%) & Male, n (\%) \\
        \hline
        ADNI-B
            & HC  & 93 & $73.26 \pm 5.33$ & 48 (51.6\%) & 45 (48.4\%) \\
            & MCI & 33 & $73.72 \pm 6.37$ & 19 (57.6\%) & 14 (42.4\%) \\
            & AD  & 26 & $75.88 \pm 7.51$ & 10 (38.5\%) & 16 (61.5\%) \\
        \hline
        ADNI-Long
            & HC  & 15 & $74.65 \pm 5.03$ & 6 (40.0\%) & 9 (60.0\%) \\
            & MCI & 23 & $74.67 \pm 6.64$ & 11 (47.8\%) & 12 (52.2\%) \\
            & AD  & 10 & $73.87 \pm 6.85$ & 4 (40.0\%) & 6 (60.0\%) \\
        \hline
    \end{tabular}
\end{table}

\subsection{Upstream Neuroimaging Preprocessing}
\label{sec:neuroimaging_preprocessing}

The neuroimaging data were preprocessed before they were provided.
This section describes that upstream preparation, not preprocessing
performed in this thesis.

Functional MRI preprocessing used the standard volume-based CONN
pipeline. It included realignment and motion correction, slice-timing
correction, and outlier detection with the Artifact Detection Toolbox
(ART). Functional images were coregistered to each participant's
high-resolution T1-weighted image and normalised to MNI-152 space with
2~mm isotropic resampling. They were then smoothed with an 8~mm FWHM
Gaussian kernel~\cite{AquileLlorensMartinezMolina}.

Denoising used anatomical CompCor (aCompCor)~\cite{Behzadi2007}. Five
principal components from white matter and five from cerebrospinal
fluid were regressed from the functional data. Nuisance regressors also
included six realignment parameters, their first-order temporal
derivatives, and ART-identified artifacts. The denoised BOLD time
series were band-pass filtered between 0.008 and
0.08~Hz~\cite{AquileLlorensMartinezMolina}.

ADNI-B and ADNI-Long underwent the same procedure but were processed
independently. Section~\ref{sec:dataset_consistency} examines the
effect on the parcellated BOLD signals.


\subsection{Brain Parcellation and BOLD Representation}
\label{sec:bold_representation}

During upstream preparation, the functional data were parcellated with
the 400-region Schaefer cortical atlas~\cite{Schaefer2018}. The
provided data therefore contain regional BOLD time series rather than
voxel-level fMRI volumes.

Each parcel is represented by one BOLD time series. An rs-fMRI sample
is represented as

\begin{equation}
    \mathbf{X} \in \mathbb{R}^{R \times T},
\end{equation}

where $R=400$ is the number of cortical regions and $T=197$ is the
number of timepoints. Each element $x_{r,t}$ is the BOLD signal for
region $r$ at timepoint $t$. Each session is thus a
$400 \times 197$ matrix.

% The 400 parcels can also be assigned to the seven large-scale networks
% described by Yeo et al.~\cite{Yeo2011}. These assignments are not model
% inputs. They are used later to interpret learned representations in
% terms of large-scale brain organisation.

Figure~\ref{fig:bold_example} shows an example of the parcellated BOLD
representation used as input to the models.

\begin{figure}[ht]
    \centering
    \includegraphics[width=0.85\textwidth]
    {example_bold_heatmap.png}
    \caption{Example parcellated rs-fMRI sample used as input to the
    models. Each row represents one of the 400 Schaefer cortical
    parcels and each column represents one of the 197 available
    timepoints.}
    \label{fig:bold_example}
\end{figure}


\subsection{Amyloid-$\beta$ and Tau Data}
\label{sec:biomarker_data}

The A$\beta$ and tau PET data were also provided after preprocessing
and parcellation. The following PET steps were performed upstream.

PET images were coregistered to the corresponding T1-weighted image and
smoothed to 8~mm FWHM. Partial-volume correction used the
Müller--Gärtner method in PETPVC~\cite{Thomas2016}.

A$\beta$ measurements came from Florbetapir (FBP) or Florbetaben (FBB)
PET scans. Standardized Uptake Value Ratios (SUVRs) used the whole
cerebellum as the reference region and were converted to Centiloid
units with tracer-specific calibration equations. Tau measurements came
from Flortaucipir PET scans. Their SUVRs used inferior cerebellar grey
matter as the reference region~\cite{AquileLlorensMartinezMolina}.

The upstream procedure parcellated PET data with the Schaefer-400
atlas. Each session is associated with regional A$\beta$ and tau
vectors,

\begin{equation}
    \mathbf{a} \in \mathbb{R}^{400},
    \qquad
    \boldsymbol{\tau} \in \mathbb{R}^{400},
\end{equation}

where $a_r$ and $\tau_r$ are the A$\beta$ and tau measurements for
region $r$.


\subsection{Input Normalisation}
\label{sec:input_normalisation}

Further normalisation was performed before the data were passed to the
models. BOLD time series and regional A$\beta$ and tau measurements
used different procedures.

Each regional BOLD time series was standardised across time using a
z-score,

\begin{equation}
    \tilde{x}_{r,t}
    =
    \frac{x_{r,t}-\mu_r}{\sigma_r},
\end{equation}

where $\mu_r$ and $\sigma_r$ are the temporal mean and standard
deviation of region $r$ in that session. This was done independently
for every region and session. It removes differences in signal scale
while preserving temporal fluctuations. Each normalised time series has
zero temporal mean and unit temporal standard deviation. The same
procedure was used for ADNI-B and ADNI-Long.

A$\beta$ and tau were instead normalised across subjects. For each
biomarker and region, the mean and standard deviation were calculated
from the ADNI-B training set only. A regional biomarker value $b_r$ was
then standardised as

\begin{equation}
    \tilde{b}_{r}
    =
    \frac{b_{r}-\mu_{r}^{\mathrm{train}}}
         {\sigma_{r}^{\mathrm{train}}},
\end{equation}

where $b_r$ is the A$\beta$ or tau value for region $r$.
$\mu_{r}^{\mathrm{train}}$ and $\sigma_{r}^{\mathrm{train}}$ are the
corresponding mean and standard deviation across training subjects.
This was done separately for each of the 400 regions and for each
biomarker.

The ADNI-B training statistics were used to normalise the validation
and test sets. They were also used for ADNI-Long when evaluating models
developed on ADNI-B. This prevents validation, test, and longitudinal
data from influencing model development.


\subsection{Consistency Between ADNI-B and ADNI-Long}
\label{sec:dataset_consistency}

Although ADNI-B and ADNI-Long underwent the same upstream neuroimaging
preprocessing procedure, the datasets were processed independently.
Consequently, sessions corresponding to the same participant and
acquisition are not numerically identical between the two datasets.
Since models developed using ADNI-B are subsequently evaluated using
longitudinal observations from ADNI-Long, the consistency between the
independently processed datasets was assessed before longitudinal
evaluation.

Sessions corresponding to the same participant and acquisition were
identified in both datasets and compared at three levels: regional BOLD
activity, static functional connectivity (FC), and sliding-window
functional connectivity dynamics (SwFCD). These comparisons were used
to determine whether differences introduced by the independent
processing affected the regional signals themselves, their static
functional relationships, or their temporal connectivity structure.

For each matched session, the normalised BOLD matrices from ADNI-B and
ADNI-Long were flattened into vectors containing all cortical regions
and timepoints,

\begin{equation}
    \mathbf{x}^{B}
    =
    \operatorname{vec}(\mathbf{X}^{B}),
    \qquad
    \mathbf{x}^{Long}
    =
    \operatorname{vec}(\mathbf{X}^{Long}),
\end{equation}

where
$\mathbf{x}^{B},\mathbf{x}^{Long}\in\mathbb{R}^{RT}$,
with $R=400$ regions and $T=197$ timepoints. Each vector therefore
contains $78{,}800$ corresponding normalised BOLD values. Similarity
between the two BOLD representations was calculated using Pearson
correlation,

\begin{equation}
    \rho_{\mathrm{BOLD}}
    =
    \operatorname{corr}
    \left(
        \mathbf{x}^{B},
        \mathbf{x}^{Long}
    \right).
\end{equation}

Static FC was calculated from the Pearson correlation between every
pair of regional BOLD time series. The unique off-diagonal elements of
the resulting FC matrices were vectorised, and similarity between the
two datasets was calculated as

\begin{equation}
    \rho_{\mathrm{FC}}
    =
    \operatorname{corr}
    \left(
        \operatorname{vec}_{u}(\mathbf{FC}^{B}),
        \operatorname{vec}_{u}(\mathbf{FC}^{Long})
    \right),
\end{equation}

where $\operatorname{vec}_{u}$ denotes vectorisation of the upper
triangular elements of a matrix, excluding the diagonal.

SwFCD was compared in the same manner using the unique off-diagonal
elements of the corresponding SwFCD matrices,

\begin{equation}
    \rho_{\mathrm{SwFCD}}
    =
    \operatorname{corr}
    \left(
        \operatorname{vec}_{u}(\mathbf{SwFCD}^{B}),
        \operatorname{vec}_{u}(\mathbf{SwFCD}^{Long})
    \right).
\end{equation}

In addition to Pearson correlation, mean absolute error (MAE) was used
to quantify the numerical difference between corresponding values. For
two vectorised representations $\mathbf{u}$ and $\mathbf{v}$ containing
$N$ corresponding elements, MAE was calculated as

\begin{equation}
    \mathrm{MAE}(\mathbf{u},\mathbf{v})
    =
    \frac{1}{N}
    \sum_{i=1}^{N}
    \left|u_i-v_i\right|.
\end{equation}

For BOLD, MAE was therefore calculated across the $78{,}800$
corresponding region--timepoint values. For FC and SwFCD, it was
calculated across the corresponding vectorised upper-triangular
elements. Pearson correlation and MAE provide complementary measures:
Pearson correlation evaluates similarity in the relative structure of
the representations, whereas MAE measures the magnitude of their
pointwise numerical differences.

Figure~\ref{fig:dataset_consistency_example} shows a representative
comparison of the BOLD, FC, and SwFCD representations from a session
present in both datasets.

\begin{figure}[ht]
    \centering
    \includegraphics[width=\textwidth]
    {b2_long_comparison.png}
    \caption{Representative comparison of a session present in both
    ADNI-B and ADNI-Long. The columns show the corresponding normalised
    regional BOLD signals, static functional connectivity (FC), and
    sliding-window functional connectivity dynamics (SwFCD). The upper
    row represents ADNI-B and the lower row represents ADNI-Long.}
    \label{fig:dataset_consistency_example}
\end{figure}

A total of eight sessions could be matched between ADNI-B and
ADNI-Long. The same BOLD, FC, and SwFCD comparisons were performed for
each shared session to assess whether the observed consistency extended
beyond an individual example. The resulting Pearson correlations and
MAE values are shown in
Figure~\ref{fig:dataset_consistency_population}.

\begin{figure}[ht]
    \centering
    \includegraphics[width=\textwidth]
    {dataset_consistency_population.png}
    \caption{Consistency between the eight sessions represented in both
    ADNI-B and ADNI-Long. Boxplots show the Pearson correlation and mean
    absolute error (MAE) between corresponding BOLD, FC, and SwFCD
    representations. Individual points represent shared sessions, and
    annotations indicate the median across sessions.}
    \label{fig:dataset_consistency_population}
\end{figure}

Across the eight shared sessions, the mean Pearson correlations were
0.770 for BOLD, 0.736 for FC, and 0.490 for SwFCD. The corresponding
mean MAE values were 0.522, 0.126, and 0.047, respectively. BOLD and FC
showed reasonably consistent correlations across the shared sessions,
indicating that their overall signal and static connectivity structure
was largely preserved between the independently processed datasets.

SwFCD showed lower and more variable Pearson correlations, indicating
that its relative structure was less consistently preserved. Within
the SwFCD representation, however, the corresponding absolute
differences remained relatively small, with a mean MAE of 0.047. This
difference between the Pearson correlation and MAE results demonstrates
that small pointwise numerical differences can coexist with weaker
preservation of the relative structure of a representation.

Overall, the independently processed datasets should not be considered
numerically equivalent. Nevertheless, the consistency check indicates
that the BOLD and static FC representations retained a reasonable degree
of similarity across the shared sessions, while larger differences were
observed in the relative structure of SwFCD. Since only eight sessions
were available in both datasets, these results are interpreted as a
dataset consistency check rather than as evidence of equivalence between
ADNI-B and ADNI-Long.

\subsection{Longitudinal Dataset Construction}
\label{sec:longitudinal_construction}

After the selection described in Section~\ref{sec:sample_selection},
the remaining ADNI-Long sessions were organised into longitudinal
pairs. Each pair contains a source and target observation from the same
participant at different visits,

\begin{equation}
    \left(
        \mathbf{X}_{t_1},
        \mathbf{X}_{t_2}
    \right),
    \qquad t_1 < t_2,
\end{equation}

where $\mathbf{X}_{t_1}$ is the source state and
$\mathbf{X}_{t_2}$ is the target state. The source session is the
earlier visit and the target session is the later visit. Their clinical
labels are referred to as the source diagnosis and target diagnosis,
respectively.

Only forward transitions between different clinical groups were kept:
HC~$\rightarrow$~MCI, HC~$\rightarrow$~AD, and
MCI~$\rightarrow$~AD. Stable and reverse transitions were excluded.

For participants with more than two sessions, all temporally ordered
source--target pairs representing valid forward transitions were
retained. A session could therefore contribute to more than one pair.
For example, for a participant with three chronologically ordered
diagnoses

\begin{equation}
    \left[
        \mathrm{MCI}_{t_1},
        \mathrm{AD}_{t_2},
        \mathrm{AD}_{t_3}
    \right],
    \qquad t_1 < t_2 < t_3,
\end{equation}

the resulting source--target pairs are

\begin{equation}
    \left(
        \mathrm{MCI}_{t_1},
        \mathrm{AD}_{t_2}
    \right)
    \quad \text{and} \quad
    \left(
        \mathrm{MCI}_{t_1},
        \mathrm{AD}_{t_3}
    \right).
\end{equation}

The pair $\left(\mathrm{AD}_{t_2},\mathrm{AD}_{t_3}\right)$ is excluded
because the diagnosis does not change. More generally, every temporally ordered source--target combination
representing one of the valid forward transitions
HC~$\rightarrow$~MCI, HC~$\rightarrow$~AD, or
MCI~$\rightarrow$~AD is retained. This allows several
progression intervals from one participant while excluding stable and
reverse transitions.

The selected ADNI-Long cohort contains 22 participants and 48 sessions,
which produced 27 progression pairs. Participants and sessions can
contribute to more than one pair or transition type.
Table~\ref{tab:longitudinal_transitions} summarises these pairs.

\begin{table}[ht]
    \centering
    \caption{Clinical progression pairs represented in the selected
    ADNI-Long cohort. The number of unique subjects within each
    transition type is not additive because the same participant may
    contribute to more than one type of clinical transition.}
    \label{tab:longitudinal_transitions}
    \begin{tabular}{lcc}
        \hline
        Transition & Unique subjects & Pairs \\
        \hline
        HC $\rightarrow$ MCI & 13 & 15 \\
        HC $\rightarrow$ AD  & 1  & 1  \\
        MCI $\rightarrow$ AD & 10 & 11 \\
        \hline
        Total & 22 & 27 \\
        \hline
    \end{tabular}
\end{table}

The interval between the source and target sessions was recorded for
each pair to describe the period over which the transition was observed.
Across the 27 longitudinal pairs, the mean interval was $2.90 \pm 1.05$ years.

As described in Section~\ref{sec:sample_selection}, participants shared
by ADNI-Long and ADNI-B were restricted to the ADNI-B test set. Thus,
longitudinal participants did not contribute to training or validation.
